\section{CHEETAH Formulation for HW/SW Co-Design}
\label{sec:methods}
CHEETAH minimizes input workload inference time under a declared pre-RTL hardware budget.
The input is a graph of logical invocations, a finite implementation library,
and a numerical contract. Hardware is shared across every invocation, operation
choices remain independent. The formulation includes the supplied schedule model and
causal execution alternatives as whole-program branches, preserving existing
feasible points while admitting richer implementations.

\subsection{Pruning, Variables, and the Encoded Features}
\label{sec:v1}\label{sec:master_lp}\label{sec:v3}
\paragraph{Pareto pruning pre-pass.}
For each operation, we discard infeasible tiles and configurations dominated
in every relevant cost, traffic, workspace, and interface dimension.
Dominance is tested within the same implementation contract and hardware class:
a lower isolated latency alone is insufficient. Timeloop or a native analytical
recipe supplies the retained coefficients once; the solver then selects among
them without invoking the mapper at each iteration. Reuse across hardware
settings requires that the same dominance relation remains valid.

Let $i\in\mathcal N$ index original invocations, $j\in\mathcal J$ storage
objects, $e\in\{0,\ldots,E-1\}$ events, and $\tau$ memory tiers. An implementation
$k$ covers members $S_k$ and specifies inputs, outputs, workspace $W_{k\tau}$,
and resource demands. $\mathcal K_P$ contains primary implementations and
$\mathcal K_R$ permitted recomputations. The latter are charged executions,
not free restoration. Table~\ref{tab:var_digest} defines the variables.

\begin{table}[t]
\centering\small
\caption{Master variables; selection and state masses are relaxed to $[0,1]$.}
\label{tab:var_digest}\label{tab:v1_scale}
\begin{tabularx}{\linewidth}{@{}l l X@{}}
\toprule
Variable & Indices & Meaning \\
\midrule
$h$ & shared chip & Compute/vector resources, memory capacities and service rates \\
$\lambda_b,z_b$ & branch $b$ & Whole-program selector and scaled branch coordinates \\
$y_k$ & implementation $k$ & Tile, mapping, batch or fusion selection \\
$s_{ke},w_{ke},f_{ke},a_{ke}$ & implementation, event & Start, work, finish and active mass \\
$D_{ie},\beta_{je}$ & invocation/object, event & Logical completion and generated data \\
$p_{j\tau e},v_{j\tau e}$ & object, tier, event & Resident state and eviction \\
$q_{j,\tau\to\tau',e}$ & object, link, event & Paid transfer and destination reservation \\
$Q^r_e,u_r,Z^r_e,d_e$ & resource, event & Demand, inverse rate, service and duration \\
$T_i,p_{et},p^W_{it},r_{it},w_{ick}$ & retained schedule block & Latency, residency, recomputation and tile--residency products \\
\bottomrule
\end{tabularx}
\end{table}

We use the same selection, residency, and latency structure throughout.
For each branch $b$, $z_b$ holds its complete coordinates and $c_b^Tz_b$ its
latency in common units. A branch weight scales every original right-hand side
and finite bound. In the causal rows below, branch indices are suppressed and
all quantities are scaled coordinates; $C^\lambda_\tau(h_b)$ is the affine
capacity with its constant multiplied by $\lambda_b$. $\mathcal R_k$ lists
required object/tier pairs, $\mathcal D_k$ completed controls, and
$\mathsf D^r_e$ is the linear demand of work and transfers. The master is given in Eq.~\eqref{eq:master}.
\begin{figure}[t]
\begingroup\small
\begin{equation}
\begin{aligned}
\min\quad &\sum_{b\in\mathcal B}c_b^Tz_b \\
\text{s.t.}\quad
 &\lambda_b\ge0,\quad\sum_b\lambda_b=1,\qquad h=\sum_b H_bz_b,\\
 &\sum_{k\in\mathcal K_P:i\in S_k} y_k=\lambda_b,\qquad
   \sum_{k\in\mathcal K_R:i\in S_k} y_k\le R_b\lambda_b,\\
 &\sum_e s_{ke}=\sum_e w_{ke}=\sum_e f_{ke}=y_k,\\
 &a_{ke}=a_{k,e-1}+s_{ke}-f_{k,e-1},\qquad 0\le w_{ke},f_{ke}\le a_{ke}\le y_k,\\
 &a_{ke}\le p_{j\tau e}-v_{j\tau e}\ ((j,\tau)\in\mathcal R_k),\qquad
   a_{ke}\le D_{ie}\ (i\in\mathcal D_k),\\
 &D_{i,e+1}=D_{ie}+\sum_{k\in\mathcal K_P:i\in S_k}f_{ke},\quad D_{i0}=0,\ D_{iE}=\lambda_b,\\
 &\beta_{j,e+1}=\beta_{je}+\sum_{k\in\mathcal K_P:j\in O_k}f_{ke},\\
 &p_{j\tau,e+1}=p_{j\tau e}-v_{j\tau e}
       +\sum_{\sigma:\sigma\to\tau}q_{j,\sigma\to\tau,e}
       +\sum_{k:j\in O_k,\delta(k)=\tau}f_{ke},\\
 &0\le v_{j\tau e}\le p_{j\tau e}\le\beta_{je},\qquad
   0\le q_{j,\tau\to\tau',e}\le p_{j\tau e}-v_{j\tau e},\\
 &\sum_j m_j\!\left(p_{j\tau e}-v_{j\tau e}
       +\sum_{\sigma:\sigma\to\tau}q_{j,\sigma\to\tau,e}\right)
       +\sum_k W_{k\tau}a_{ke}\le C^\lambda_\tau(h_b),\\
 &Q^r_e=\mathsf D^r_e(q,w),\quad (h_{br},u_r)\in\mathcal R_b^{\lambda_b},\quad
      (Q^r_e,u_r,Z^r_e)\in\mathcal M_{br}^{\lambda_b},\\
 &\sum_{r\in\mathcal P_a}Z^r_e\le d_e\quad(\text{each shared service pool }a),
       \qquad F_bh_b\le g_b\lambda_b,\\
 &B_bz_b=b_b\lambda_b,\quad J_bz_b\le j_b\lambda_b
       \quad(\text{boundary, backing and implementation rows}),\\
 &A_b^Pz_b=a_b^P\lambda_b,\quad G_b^Pz_b\le g_b^P\lambda_b
       \quad(\text{retained schedule branches}),\\
 &\ell_{bj}\lambda_b\le z_{bj}\le u_{bj}\lambda_b
       \quad(\text{each finite original bound}).
\end{aligned}
\tag{Master}\label{eq:master}
\end{equation}
\stepcounter{equation}
\endgroup
\end{figure}
Causal branches use the cover, lifecycle and service rows; retained schedule
branches use their complete original rows $A_b^P,G_b^P$. These are alternatives
inside one program, not extra causal restrictions on the old schedule model.
Appendix~\ref{app:master_rows} expands the abbreviated rows, including the
original joint tiling--placement block. Setting one $\lambda_b=1$ embeds every
point of that branch with unchanged hardware and objective. A fractional
mixture is a relaxation, not a multi-chip execution; recovery selects one branch.

\subsection{Features and Their Mathematical Encoding}
\label{sec:v3_levers}\label{sec:v0}
\paragraph{Joint tiling and fusion.}
For an isolated operation, the cover row reduces to $\sum_c y_{i,c}=1$.
A tile's traffic savings depend on residency. In the retained schedule block,
$w_{ick}$ represents $y_{ic}p_{\mathrm{assoc}(i,k),o(i)}$ through
\begin{equation}
0\le w_{ick},\quad w_{ick}\le y_{ic},\quad w_{ick}\le p,\quad
w_{ick}\ge y_{ic}+p-1.
\label{eq:mccormick}
\end{equation}
These McCormick inequalities give the convex hull of one bounded bilinear graph
and are exact at integral factors~\citep{mccormick1976computability}; they do not
make every fractional solution executable. Tile-dependent footprints use the
same construction. Native fused atoms instead carry their complete input,
output, workspace and service recipe, so fusion cannot silently remove work.

\paragraph{Branch and bound.}
The envelopes above are exact only at integral factors, so a point the
relaxation returns need not be a design that can be built. A fractional tile
selection corresponds to no single tiling; a fractional menu weight mixes two
memory technologies or array aspects that no die realizes; and a fractional
branch weight is a blend of execution alternatives rather than one executable
schedule. In other words, the relaxation bounds the achievable latency but does
not by itself exhibit a valid physical chip. We therefore close the gap with a
complete branch-and-bound search over the discrete part of the design, which is
what lets us claim an actual physically feasible chip: the one-hot menus and
tile selectors are branched, each node re-solving the same continuous master
with those selectors pinned, which keeps every node a linear program while the
branching restores integrality. The relaxation supplies the lower bound at each
node and any integral point that survives recovery and replay supplies the upper
bound, so the search terminates with a design that is both buildable and
bracketed. Branching is applied to the discrete selectors only; the shared
hardware sizing below stays continuous wherever the underlying physics is
continuous.

\paragraph{Shared hardware and resource epigraphs.}
Compute, vector lanes, SRAM, and delivered memory service are shared decisions
under linear area/power proxies $F_bh_b\le g_b$. Finite implementation or
technology menus use disaggregated copies with linking equalities and scaled
bounds. Resource epigraphs charge the maximum of concurrent service demands,
while demands on the same port add. A per-implementation epigraph retains the
compute/memory trade-off before summation; arbitrary overlap is not inferred
from fractional sharing. The native replay uses its declared serial kernel and
copy semantics.
\label{eq:v3_caps}\label{eq:v3_epigraph}

\paragraph{Reciprocal and product envelopes.}
For a rate $r\in[\ell,h]$, $\mathcal R$ bounds $u\approx1/r$ by tangents
$u\ge2/a-r/a^2$ and the global secant $u\le(\ell+h-r)/(\ell h)$.
$\mathcal M$ bounds $Z\approx Qu$ by the four McCormick inequalities on
$[0,\bar Q]\times[1/h,1/\ell]$. All coefficients are constants, so the program
is linear. Outward coefficient rounding preserves graph containment.
These are relaxations of physical service; replay reinstates the exact rates
and products. Discrete bandwidth-menu copies are exact when their selector is
integral.

\paragraph{Time-indexed placement and KV across steps.}
Residency $p_{j\tau e}$ changes with events rather than being fixed for the whole
run. KV production creates immutable append objects; each decode attention
operation depends on its required prefix. The state rows preserve earlier
entries while the prefix grows. Capacity accounts for resident data, incoming
copies, and active workspace simultaneously. Long skip connections and
multi-fanout tensors use the same object identities and consumer readiness.

\paragraph{Spill, restoration, and rematerialization.}
Copies retain a valid source, reserve their destination, consume service, and
become readable only after completion. Boundary rows require the declared
terminal backing. A permitted repeat executes the original producer again and
pays its work, input and workspace demands; it restores storage without
advancing primary logical completion twice. Keeping full KV history in a
slower tier is therefore different from deleting history. The native results
credit paid KV restoration and retention, not a lossy KV budget.

\paragraph{Batching and causal scheduling.}
A batch column covers several compatible invocations and loads shared parameters
once where its recipe permits. Cover constraints retain every member;
start/finish and readiness rows prevent using future tokens or expert selections
before their router dependencies complete. Event placement can vary within the
finite library. The current native execution uses a recovered serial schedule;
a larger scheduling domain does not imply that its optimum has been found.

\paragraph{Implementation menus and recovery.}
Precision, sparsity, array aspect, memory technology, and validated overlap
policies enter as implementation choices with their own numerical and resource
contracts. Accuracy-changing cache eviction or approximate reuse needs separate
quality evidence; it is pinned off in the full-input native result. Larger
cross-GEMM fusion and unrestricted overlap have no credited gain there.
Recovery pins hardware and mappings, schedules transfers, and replays the
original obligations. The master LP itself is continuous, and we run a complete
branch-and-bound search over its discrete selectors to recover designs that are
actual physically feasible chips: each node re-solves the continuous master with
the one-hot menu and tile selectors pinned, the relaxation supplies the node
bound, and a recovered point that survives replay supplies the incumbent, so the
search closes with a buildable design and a certified bracket.

\subsection{LLM Architect Diagnostic Outer Search}
\label{sec:llm_outer}
The LLM architect reads a ledger of resource prices, slacks, invalid trials,
and replay outcomes, then proposes candidate implementations or search
neighborhoods. Deterministic validation checks each proposal before evaluation.
In the sequential engineering pipelines, tools such as Vizier~\citep{google_vizier} searches those neighborhoods through Bayesian search of thousands of trials.
the master also represents hardware choices directly. Dual prices guide search,
but an LLM proposal is not a complete pricing oracle: a full-domain optimality
claim requires checking all omitted alternatives. The original controller and
prompts are documented in Appendix~\ref{app:llm_loop}.
